\chapter{Outils numériques et synthèse}
\label{exercices:20}
\section*{Exercice 20.1 -- Tableur et recopie}
Les valeurs sont en \texttt{A2:A6}, leurs effectifs en \texttt{B2:B6}. Écrire les formules pour le total en \texttt{B8}, la fréquence en \texttt{C2} à recopier et la moyenne pondérée. Expliquer les dollars.

\section*{Exercice 20.2 -- Construction dynamique}
Décrire une construction dynamique d'une configuration de Thalès dont les hypothèses restent vraies quand on déplace les sommets. Les rapports affichés égaux prouvent-ils le théorème ?

\section*{Exercice 20.3 -- Un modèle de remplissage}
Un cylindre de rayon intérieur $0{,}5\unite m$ et de hauteur $2\unite m$ est rempli à $20\unite{L/min}$. Calculer sa capacité exacte, le temps pour le remplir à $60\,\%$ et la fonction donnant le volume stocké avant saturation.

\section*{Exercice 20.4 -- Aire maximale, complément}
Un rectangle a un périmètre de $40\unite{cm}$. Montrer que son aire est $A(x)=100-(x-10)^2$, puis en déduire sa plus grande aire sans utiliser de dérivée.
